\section{Implementation-Level Validation: Cycle-Accurate RTL}
\label{sec:rtl_validation}

The results so far certify \emph{model-level} quantities computed inside one
cost-model convention; as we note in \cref{app:mfu}, a real measured
utilization would require an executable implementation of the selected
hardware/schedule pair. This section supplies that check. We take the
committed schedule and the datapath geometry the LP sizes, express them as
\emph{synthesizable} SystemVerilog, simulate them cycle-accurately with
Verilator~5.050~\citep{snyder_verilator}, and compare against a reference model
written in C\texttt{++}. We call this \emph{analytical-model-to-RTL cycle
validation}: the RTL confirms that the LP's schedule is physically
realizable, that its time-indexed memory constraints hold byte-for-byte in
hardware, and that the utilization \emph{mechanisms} the LP charges are the
ones a real datapath exhibits. We do \emph{not} treat RTL cycle counts as
absolute ground truth (the analytical per-layer cache floors are uniformly
conservative by $192\times$ to $2475\times$, quantified below), and we make
no tape-out claim. \Cref{sec:asap7} then closes the remaining gap on the
cost-model side, taking the sized datapath through logic synthesis and
placement at a $7$~nm predictive PDK to test the area and power coefficients
themselves.

\paragraph{Experiment 1: schedule tape replay (realizability and memory
legality).}
We compile the full committed DeepSeek-V3~\citep{deepseekv3report} schedule
at a TPU-v3-class configuration ($\CGM = 8$~MiB) into a $1{,}383$-instruction
event tape over the vocabulary \{\textsc{execute}, \textsc{load},
\textsc{free}, \textsc{evict}, \textsc{reload}, \textsc{rematerialize}\}, and
execute it on a cycle-accurate SystemVerilog event engine (a Global-Buffer
allocator, a bandwidth-limited DMA model, and a compute model) under
Verilator. Each instruction is cross-checked against the analytical residency
model layer-by-layer, and a deliberately over-committed tape confirms the
capacity checker is live (it flags exactly one violation). Three pre-stated
criteria pass (\cref{tab:rtl_validation}), and \cref{tab:rtl_levers} reports, decision by
decision, that every planned event executed. \emph{Memory legality:} the tape
runs to completion with \textbf{zero} capacity violations, and the RTL peak
on-chip occupancy equals the analytical prediction \emph{exactly}
($7{,}602{,}176$~B $= 7.25/8$~MiB, $9.4\%$ headroom). \emph{Cycles:} the total
RTL cycle count lands within $+0.000001\%$ of the rounded LP objective, the
entire gap being the $16{,}431$ decoder/DMA sequencing cycles ($10^{-8}$ of
the total); i.e., the rounded schedule realizes with \emph{no hidden
serialization or stall} beyond the model's accounting. \emph{Two-tape ratio:}
against a stream-everything realization of the same operation order, the
RTL-measured latency ratio is $1.0723\times$, matching the analytical ratio to
$8$ significant digits. This ratio is scoped to the pin/reload channel alone
(both tapes hold the per-operation-optimal tiling fixed, so joint tile
selection, the dominant channel elsewhere, is deliberately held equal).
Together these establish that the LP's time-indexed residency model
(\cref{sec:v0}) is not merely roofline-feasible but executes byte- and
cycle-faithfully on hardware that enforces the capacity constraint at every
step.

\begin{table}[t]
\centering
\caption{Did each decision fire? Per-decision plan-versus-execution counters from the DeepSeek-V3 tape replay ($\CGM = 8$\,MiB). Every event the LP plans is executed by the RTL engine exactly once, the capacity row is checked in hardware at every step, and the decisions the LP chose not to use at this cell are confirmed absent while their opcodes are proven live on a synthetic tape. The bound says how good the plan is; the counters say the plan is what ran.}
\label{tab:rtl_levers}
\scriptsize
\setlength{\tabcolsep}{5pt}
\renewcommand{\arraystretch}{1.22}
\definecolor{bandrtl}{RGB}{230,239,253}
\definecolor{rowalt}{gray}{0.955}
\begin{tabular}{@{}l l l l@{}}
\toprule
\rowcolor{bandrtl}
\textbf{Decision} & \textbf{LP plan} & \textbf{RTL executed} & \textbf{Check} \\
\midrule
Tile choice ($y$)            & $903$ \textsc{execute}               & $903$                 & \cmark~ids match $903/903$ \\
\rowcolor{rowalt}
Residency ($p$, $p^W$)       & $180$ \textsc{load}, $299$ \textsc{free}, $2{,}709$ pins & $180$, $299$ & \cmark~$2{,}709/2{,}709$ pins \\
Fusion (operand kept resident) & $418$ of $903$ ops                 & $418$                 & \cmark \\
\rowcolor{rowalt}
Eviction / reload            & $0$ / $0$                            & $0$ / $0$             & \cmark~opcodes live \\
Rematerialization ($r$)      & $0$                                  & $0$                   & \cmark~opcode live \\
\rowcolor{rowalt}
Capacity row                 & $1{,}083$ hardware checks            & $0$ violations        & \cmark~injected overflow flags $1$ \\
Peak occupancy               & $7{,}602{,}176$\,B                   & $7{,}602{,}176$\,B    & \cmark~byte-exact, trace within $4$\,B \\
\rowcolor{rowalt}
Total cycles                 & $1.6437\times10^{12}$                & $+16{,}431$ cycles    & \cmark~$+0.000001\%$, accounted \\
\bottomrule
\end{tabular}
\end{table}

\paragraph{Experiment 2: datapath utilization mechanism.}
We next build one weight-stationary INT8 systolic PE cluster in synthesizable
SystemVerilog ($4\times4$ MAC array, INT8$\times$INT8$\rightarrow$INT32, with
L1 buffers and a tile-loop FSM; plus a $16\times16$ variant), golden-verified
against a plain-C\texttt{++} reference convolution ($8/8$ runs bit-exact).
Simulated on EfficientNet-B7~\citep{tan2019efficientnet} layers, the RTL
reproduces the two utilization mechanisms the cost model charges for. First, the
\emph{depthwise structural ceiling}: on a broadcast systolic array a depthwise
convolution can engage only $4$ of $16$ MAC lanes, and the RTL measures a
steady-state depthwise utilization of $24.9\%$ against the $25\%$ analytic
ceiling ($6.13\%$ vs.\ $6.25\%$ on the $16\times16$ variant), while dense
$1\times1$ layers hold $97$--$99\%$ on the identical RTL. Second, the
predicted \emph{geometry contrasts}: the measured lane law reproduces the
$4.00\times$ systolic-vs-grid contrast to $4.057\times$, and extends the
cache's $18\times$ depthwise contrast to within $1.22\times$. Crucially, the
same runs show the analytical cache cycle \emph{floors} to be $192\times$ to
$2475\times$ conservative (the RTL executes the identical work that much
faster), which is precisely why RTL-derived claims are scoped to these
utilization \emph{contrasts} and to the schedule replay above, never to
absolute cycle agreement. The mechanism is what matters here: the utilization
losses the LP models (\cref{sec:methods}) are the ones the real datapath
exhibits.

\begin{table}[t]
\centering
\small
\caption{Analytical-model-to-RTL cycle validation. Synthesizable
SystemVerilog simulated cycle-accurately with Verilator~\citep{snyder_verilator}
and checked against a C\texttt{++} reference. The RTL confirms schedule
realizability, byte-exact memory legality, and the LP's utilization
mechanisms; it is \emph{not} used as an absolute cycle oracle.}
\label{tab:rtl_validation}
\begin{tabular}{@{}lll@{}}
\toprule
\textbf{Quantity (RTL vs.\ analytical)} & \textbf{Measured} & \textbf{Verdict} \\
\midrule
\multicolumn{3}{@{}l}{\emph{Exp.\ 1: DeepSeek-V3 schedule replay
($1{,}383$ instr., $\CGM=8$~MiB)}} \\
Capacity violations & $0$ & pass \\
Peak on-chip occupancy & $7{,}602{,}176$~B (exact) & byte-exact \\
Total cycles vs.\ rounded LP obj. & $+0.000001\%$ & no hidden serialization \\
Two-tape ratio (pin/reload only) & $1.0723\times$ ($8$-digit match) & pass \\
\midrule
\multicolumn{3}{@{}l}{\emph{Exp.\ 2: EfficientNet-B7 INT8 systolic cluster
(golden $8/8$)}} \\
Depthwise MAC-lane utilization & $24.9\%$ vs.\ $25\%$ ceiling & mechanism confirmed \\
Systolic-vs-grid geometry contrast & $4.057\times$ vs.\ $4.00\times$ & mechanism confirmed \\
Absolute cache cycle floors & $192$--$2475\times$ conservative & scope guard \\
\bottomrule
\end{tabular}
\end{table}

\subsection{Physical-design validation: post-synthesis area and the linearity
of the cost model}
\label{sec:asap7}

Experiments 1 and 2 validate the \emph{schedule} the program commits to. A
second, independent question is whether the \emph{cost model} that evaluates the
hardware is itself faithful: the program sizes a chip through the coefficients
$k_{\mathrm{MAC}}$ (area per MAC) and $\pi_{\mathrm{MAC}}$ (power per MAC), and
every manufacturability gate and the area and power caps are expressed in them.
If those coefficients are wrong, the certified designs are certified against the
wrong physics. We therefore take the datapath the program sizes, write it as
synthesizable RTL, and push it through a full open-source RTL-to-placement flow
at an advanced predictive node.

\paragraph{Setup.}
We implement a parameterized weight-stationary INT8 systolic tile (an
$N \times N$ array of processing elements, each holding a resident $8$-bit
weight, streaming $8$-bit activations east and $32$-bit partial sums south, one
MAC per PE per cycle in steady state) and synthesize, floorplan, place, and
timing-repair it with OpenROAD~\citep{ajayi2019openroad} against the ASAP7
$7$~nm predictive PDK~\citep{clark2016asap7} at a $1$~GHz target, for
$N \in \{16, 32, 64\}$, that is $256$ to $4{,}096$ MACs. We report post-detailed-placement
extracted areas. We do not run clock-tree synthesis or routing, and we make no
tape-out claim: the quantity under test is the area coefficient and its scaling
law, both of which are fixed at placement.

\paragraph{Result 1: the linear area model is structurally correct.}
\Cref{tab:asap7} reports extracted cell area at the three sizes. A linear fit
gives $A_{\mathrm{cell}} = 82.08\,\mathrm{MACs} - 206.5\ \mu\mathrm{m}^2$ with
$R^2 = 0.999999$, and the extracted area \emph{per MAC} varies by only $0.46\%$
across a $16\times$ range in MAC count. This is a stronger statement than
agreement on a single number: the program assumes area is
\emph{linear} in the MAC count, with the same slope at every chip size, and
that assumption is exactly what a placed datapath exhibits. The linear area row
in the program is therefore not a convenient approximation but a measured
property of the datapath it sizes.

\paragraph{Result 2: the area coefficient is anchored and conservative.}
The fitted slope is $8.21 \times 10^{-5}\ \mathrm{mm}^2$/MAC of standard-cell
area. Placed designs occupy a core larger than their cell area; our floorplans
sit at $41\%$ placement utilization, giving $2.0 \times 10^{-4}\
\mathrm{mm}^2$/MAC of core area. The value the program uses,
$k_{\mathrm{MAC}} = 1.5 \times 10^{-4}\ \mathrm{mm}^2$/MAC, lies between these
two and corresponds to the cell area at $54.7\%$ placement utilization, a
typical density for a datapath of this kind. The coefficient is thus bracketed
by extracted silicon area from below and by our own conservative floorplan from
above.

\paragraph{Result 3: the power coefficient, and why we anchor it to silicon
rather than to this flow.}
Power at the same operating point is instructive in a different way. Reporting
placement-level power on this tile with flat random switching activity yields
$2.0$ to $6.5$ pJ per MAC operation as the assumed input activity rises from
$5\%$ to $50\%$, against the $0.40$ pJ/MAC-op implied by
$\pi_{\mathrm{MAC}}$. Rather than adjudicate this within the flow, we test it
against an external physical bound. Shipped accelerators deliver INT8
arithmetic at $0.35$ pJ/op (H100, $4$~nm), $0.64$ pJ/op (A100, $7$~nm), and
$1.27$ pJ/op (TPU-v4i, $7$~nm), and these are \emph{whole-chip} figures that
include memory, cache, and control. Our number is for the multiplier array
alone, at the same node, and exceeds an entire A100 per operation by
$3.2$--$10.1\times$; a bare array cannot consume more energy per operation than
the complete system that contains it. The measurement, not the coefficient, is
therefore out of range: with a generic synthesized multiplier, no clock gating
or operand isolation, and flat random activity in place of correlated INT8
data, this flow reports an upper bound rather than a representative figure. The
program's $\pi_{\mathrm{MAC}} = 0.40$ pJ/MAC-op sits where an array-only
coefficient should, at the same order as and slightly below whole-chip shipped
silicon, and we anchor it there~(\cref{app:coeff_provenance}).

\paragraph{Synthesis of the three hardware verification results.}
Design-space-exploration frameworks are conventionally validated in three ways:
against measured taped-out chips, against post-synthesis extraction, and
against other frameworks~\citep{mei2021zigzag}. The Timeloop-derived Pareto
data places this work in the third category and the silicon-anchored
coefficients in the first; this section supplies the second. Concretely, it
converts the area side of every gate-legality claim in this paper, and of the
certified designs in \cref{sec:experiments}, from a literature-shaped constant
into a quantity extracted from RTL-to-placement at $7$~nm, together with a
validated scaling law; and it leaves the power side anchored to shipped
silicon, with an explicit account of why this particular flow overstates it.

\begin{table}[h]
\centering
\footnotesize
\caption{Post-placement extracted area of the INT8 systolic tile at ASAP7
$7$~nm. Area per MAC is constant to within $0.46\%$ over a $16\times$ range in
MAC count; the linear fit has $R^2 = 0.999999$.}
\label{tab:asap7}
\begin{tabular}{@{}rrrrr@{}}
\toprule
$N$ & MACs & Cell area ($\mu\mathrm{m}^2$) & Cell area/MAC ($\mathrm{mm}^2$) &
Core area/MAC ($\mathrm{mm}^2$) \\
\midrule
$16$ & $256$   & $20{,}981$  & $8.196\times10^{-5}$ & $1.986\times10^{-4}$ \\
$32$ & $1{,}024$  & $83{,}629$  & $8.167\times10^{-5}$ & $1.986\times10^{-4}$ \\
$64$ & $4{,}096$  & $336{,}049$ & $8.204\times10^{-5}$ & $1.998\times10^{-4}$ \\
\midrule
\multicolumn{3}{@{}l}{Linear fit ($R^2 = 0.999999$)} &
$8.208\times10^{-5}$ & \\
\multicolumn{3}{@{}l}{Program value $k_{\mathrm{MAC}}$} &
\multicolumn{2}{l}{$1.5\times10^{-4}$ (cell area at $54.7\%$ utilization)} \\
\bottomrule
\end{tabular}
\end{table}

\subsection{Measuring the compute-energy coefficients}
\label{sec:energy_coeff}

The area result above validates $k_{\mathrm{MAC}}$. The program also charges for
compute \emph{energy}, and one property of that dual-guided scoring deserves a measurement
rather than an assumption. Reducing arithmetic precision is usually modelled as
a reduction in bytes moved, and in an earlier form of our program precision
affected memory traffic alone. That is incomplete: a narrower multiplier also
switches less, so precision should reduce compute energy as well. We add that
term, and because a multiplier's energy is often argued to follow its
partial-product count, we could have set the ratio analytically. Instead we
measured it.

\paragraph{Method.}
We synthesize and place six variants of the tile at ASAP7: three operand-width
combinations (W8A8, W4A8, W4A4) crossed with two power-management levels
(ungated, and clock-gated with operand isolation). Power is reported at a matched
input activity so that the ratios are comparable. As a check on the flow, the
W8A8 ungated variant reproduces the independently synthesized tile of
\cref{tab:asap7} to within $0.1\%$ despite a different source file, module name
and container.

\paragraph{Result: the analytic ratio is wrong in both directions.}
\Cref{tab:energy} reports the measurements. The partial-product model predicts
energy ratios of $0.50$ for W4A8 and $0.25$ for W4A4 against the W8A8 baseline.
The measured ratios are $0.376$ and $0.358$. Narrowing the \emph{weight} saves
more than the model predicts, because switching concentrates in the multiplier
tree; narrowing the \emph{activation} saves far less, because the $32$-bit
accumulator and the partial-sum pipeline registers do not narrow with the
operands, so the second halving buys only a further $5\%$. We therefore charge
compute energy with the measured ratios rather than the analytic ones. The
$\mathrm{W4A16}$ entry is not directly measured and is held at a conservative
$0.5$ rather than extrapolated in our own favour.

\paragraph{A negative result on gating, and why we do not include it.}
The same matrix was intended to characterize a clock-clock-gating decision, and it cannot. Under
a flat synthetic activity model the gated variants measure as \emph{more}
expensive than the ungated ones ($1.40\times$ the power at W8A8), which is not
physical. The cause is that a flat activity assumption charges for the gating
cells while never crediting the idle cycles they exist to skip; the benefit of
gating is visible only under a realistic idle pattern, which requires
switching activity captured from a simulation of the actual workload. The area
cost of gating is measured and real ($+22\%$ at W4A8, $+20\%$ at W4A4), but
without a credible power benefit to weigh against it we do not introduce the
decision. We report the attempt because the failure is informative: it delimits
what this class of measurement can and cannot establish.

\paragraph{Scope.}
These measurements sharpen a coefficient; they do not change any reported
speedup, and the reason is visible in the dual variables
(\cref{sec:binding}): the power constraint is slack at every design we report,
so by complementary slackness a change that relieves it cannot move the
optimum. The value of the measurement is provenance. The best measured variant
consumes $1.22$~pJ per MAC operation against the $0.40$ implied by
$\pi_{\mathrm{MAC}}$, a factor of three rather than the order of magnitude that
an ungated tile driven by random data suggested, with the residual attributable
to predictive-PDK pessimism, a generic rather than optimized multiplier, and
synthetic activity. We continue to anchor $\pi_{\mathrm{MAC}}$ to shipped
silicon.

\begin{table}[h]
\centering
\footnotesize
\caption{Measured compute energy and area of the tile at ASAP7 $7$~nm
($N=16$, $1$~GHz, matched input activity). Ratios are against the W8A8 ungated
variant. The analytic partial-product model predicts $0.50$ and $0.25$ for the
W4A8 and W4A4 energy ratios; both are wrong.}
\label{tab:energy}
\begin{tabular}{@{}lrrrr@{}}
\toprule
Variant & Power (W) & pJ/MAC-op & Energy ratio & Area ratio \\
\midrule
W8A8, ungated & $0.872$ & $3.41$ & $1.000$ & $1.001$ \\
W4A8, ungated & $0.328$ & $1.28$ & $\mathbf{0.376}$ & $0.726$ \\
W4A4, ungated & $0.312$ & $1.22$ & $\mathbf{0.358}$ & $0.625$ \\
W8A8, gated   & $1.220$ & $4.77$ & $1.399$ & -- \\
W4A8, gated   & $0.347$ & $1.36$ & $0.398$ & $0.887$ \\
W4A4, gated   & $0.312$ & $1.22$ & $0.358$ & $0.752$ \\
\bottomrule
\end{tabular}
\end{table}
